\subsection{Calendar-block uncertainty of the spatial comparison}
Using the paired block procedure in Section~\ref{sec:block_score_uncertainty}, the California OC--EP
comparison separates a weak first-fold contrast from the larger subsequent
deficits. The intervals below concern the joint score.

\begin{table}[htbp]
\centering
\begin{tabular}{lrrr}
\toprule
Fold & Contrast (nats) & Paired 95\% interval & Nominal sign-flip $p$ \\
\midrule
F0 & -6.289 & $[-21.790,3.296]$ & 0.57656 \\
F1 & -31.743 & $[-77.229,-3.991]$ & 0.02558 \\
F2 & -13.020 & $[-22.823,-4.861]$ & 0.00366 \\
\midrule
Pooled & -51.051 & $[-99.930,-17.264]$ & 0.00018 \\
\bottomrule
\end{tabular}
\caption{California OC minus EP: 90-day paired calendar-block uncertainty. The pooled bootstrap resamples within each fold.}
\label{tab:ca_oc_ep_uncertainty}
\end{table}

The pooled joint deficit is -51.051~nats; its 180- and 365-day intervals are $[-96.894,-16.999]$ and $[-93.188,-17.257]$, with sign-flip values $p=0.00130$ and $p=0.00726$. The location component separately totals -51.387~nats, with a 90-day interval $[-98.567,-18.423]$ and $p=0.00006$. Thus the roughly 51-nat joint loss is also visible after preserving within-block event dependence.
